\documentclass[11pt, oneside]{article}

\pdfminorversion=7
\usepackage{HomeworkStyle}
\hypersetup{
    pdftitle={Homework 1 for CSED490V Foundations of Continuous Computation},
    pdfauthor={Jungwoo Kim},
    pdfsubject={Homework 1: computable real numbers and functions},
    pdfcreator={LaTeX}
}

\begin{document}
\sethomeworknumber{1}
\homeworktitle{Homework 1}

When you prove the computability of a real number or a real function, you should
explicitly provide an algorithm that is required in the definition of
computability. In the algorithm, you can only use rational numbers, basic
operations over rationals (i.e., $+$, $-$, $\times$, $/$, $\code{abs}$),
comparisons between rationals (e.g., $=$, $<$, $\le$), and basic programming
constructs (i.e., \code{if}, \code{for}, \code{while}).

\begin{problem}{1}
In this problem, ``computable'' means Borel-computable.
\begin{problemsubparts}
    \item Prove that if $x, y\in\R$ are computable, then $x+y$ is computable.
    \item Prove that $\sqrt{k}$ is computable for all integers $k\ge2$.

    \begin{hint}
    Define $(x_m)_{m\in\Nzero}\subseteq\Q$ by
    $x_0\coloneqq k$ and
    \[
      x_{m+1}\coloneqq\frac12\left(x_m+\frac{k}{x_m}\right).
    \]
    Show that there exists $\alpha>0$ such that
    $0<x_m-\sqrt{k}<k\alpha^m$ for all $m\in\Nzero$.
    Use this result to prove the original claim.
    \end{hint}
    % The original indexes (x_m) by N but specifies x_0; N_0 is intended.
\end{problemsubparts}
\end{problem}

\begin{solution}
\begin{problemsubparts}
    \item Let $p, r\colon\Nzero\to\Q$ be computable sequences satisfying
    \[
        |p(n)-x|<2^{-n}, \qquad |r(n)-y|<2^{-n}
        \quad(n\in\Nzero).
    \]
    On input $n\in\Nzero$, compute $p(n+1)$ and $r(n+1)$ and return
    $q(n)\coloneqq p(n+1)+r(n+1)$. This algorithm uses only rational
    addition and calls to the given computable approximation procedures.
    Moreover,
    \[
    |q(n)-(x+y)|
      \le |p(n+1)-x|+|r(n+1)-y|
      <2^{-(n+1)}+2^{-(n+1)}=2^{-n}.
    \]
    Hence $(q(n))_{n\in\Nzero}$ computes $x+y$.

    \item Fix an integer $k\ge2$, and define $x_0=k$ and
    \[
        x_{m+1}=\frac12\left(x_m+\frac{k}{x_m}\right)
        \qquad(m\in\Nzero).
    \]
    The algorithm starts at the rational number $k$ and computes each
    successive $x_m$ by positive rational arithmetic. Put $s=\sqrt{k}$.
    Because $k>s>0$,
    \[
      x_{m+1}-s
      =\frac{x_m^2-2sx_m+s^2}{2x_m}
      =\frac{(x_m-s)^2}{2x_m}>0
    \]
    whenever $x_m>s$. By induction, $x_m>s$ for every $m$.
    Since $0<x_m-s<x_m$, this identity also gives
    \[
      0<x_{m+1}-s<\frac{x_m-s}{2}.
    \]
    Iteration yields
    \[
      0<x_m-s\le\frac{k-s}{2^m}<\frac{k}{2^m}
      \qquad(m\in\Nzero),
    \]
    where the strict inequality at $m=0$ follows from $k-s<k$.
    Thus the requested estimate holds with $\alpha=\tfrac12$.

    Given $n\in\Nzero$, find an integer $m\ge0$ by starting at $m=0$
    and incrementing $m$ while $k2^{-m}\ge2^{-n}$.
    This search terminates because repeated division by $2$ sends
    $k2^{-m}$ to $0$, and its test uses only rational comparisons.
    Return the rational number $x_m$, computed by $m$ Newton steps.
    Then
    \[
       |x_m-\sqrt{k}|<k2^{-m}<2^{-n},
    \]
    proving the Borel computability of $\sqrt{k}$.
\end{problemsubparts}
\end{solution}

\begin{problem}{2}
Prove that the following problem is decidable (i.e., there is an algorithm
solving the problem): Given computable functions $f, g\colon\N\to\Q$
representing distinct computable reals $x, y\in\R$ (i.e.,
$|f(n)-x|<2^{-n}$ and $|g(n)-y|<2^{-n}$ for all $n\in\N$,
and $x\ne y$), decide whether $x<y$ or $x>y$.
\end{problem}

\begin{solution}
For $n=1,2,3, \ldots$, compute $a\coloneqq f(n)$,
$b\coloneqq g(n)$, and $r\coloneqq2^{-n}$, and perform the following tests:
\begin{enumerate}[label=(\arabic*), leftmargin=2.5em, topsep=4pt, itemsep=2pt]
    \item If $a+r\le b-r$, return $x<y$.
    \item If $b+r\le a-r$, return $x>y$.
    \item Otherwise, increment $n$ and repeat.
\end{enumerate}
Only rational arithmetic, comparisons, and a loop are required.
If the first condition holds, the approximation guarantees imply
\[
 x<a+r\le b-r<y,
\]
and hence the answer is correct; the second condition is symmetric.

To prove termination, let $d\coloneqq|x-y|>0$.
Choose $n$ sufficiently large that $4\cdot2^{-n}<d$.
If $x<y$, then
\[
 f(n)+2^{-n}<x+2\cdot2^{-n}
   <y-2\cdot2^{-n}<g(n)-2^{-n},
\]
so the first test succeeds. If $x>y$, the symmetric argument makes the
second test succeed. Thus the algorithm always halts and returns the
correct order under the promise $x\ne y$.
\end{solution}

\begin{problem}{3}
\begin{problemsubparts}
    \item Prove that $\exp\colon\R\to\R$ is computable over the entire
    $\R$. Your proof should be self-contained: it must include all details,
    and must not refer to any part of the proof covered in class.
    \item See \code{p3.ipynb}.
\end{problemsubparts}
\end{problem}

\begin{solution}
\begin{problemsubparts}
    \item We give a rational-only oracle algorithm, including an explicit
    computable error budget, that works for every $x\in\R$.
    Let $\alpha\colon\Nzero\to\Q$ be an arbitrary rational name of $x$,
    so that $|\alpha(i)-x|<2^{-i}$ for all $i\in\Nzero$.
    Fix an output precision $n\in\Nzero$.

    First query $q_0\coloneqq\alpha(0)$. Starting at $B=3$ and
    increasing $B$ by $1$ as needed, find an integer $B$ such that
    \[
        B>|q_0|+2.
    \]
    The search uses rational comparisons and terminates. For every
    $i\in\Nzero$, the approximation inequalities give
    \[
       |x|<|q_0|+1<B,
       \qquad |\alpha(i)|<|x|+1<B.
    \]
    Set $C\coloneqq3^B\in\N$, computing it by $B$ successive
    multiplications by $3$.
    To justify this bound without using a computability result for
    $\exp$, observe that for $j\ge2$, $j!\ge2^{j-1}$, with a strict
    inequality for $j\ge3$. Therefore
    \[
      e=\sum_{j=0}^{\infty}\frac1{j!}
      <2+\sum_{j=2}^{\infty}2^{-(j-1)}=3.
    \]
    Since $\exp$ is increasing and $\exp(B)=e^B$ for integer $B$,
    \[
      0<\exp(t)\le\exp(B)=e^B<C
      \qquad(-B\le t\le B).
    \]
    In particular, the mean value theorem shows that, for
    $u, v\in[-B, B]$,
    \[
       |\exp(u)-\exp(v)|\le C|u-v|.
       \tag{1}
    \]
    Taylor's theorem with the Lagrange remainder, applied to
    $P_d(t)\coloneqq\sum_{j=0}^{d}t^j/j!$, similarly gives
    \[
       |\exp(t)-P_d(t)|
       \le C\frac{B^{d+1}}{(d+1)!}
       \qquad(|t|<B).
       \tag{2}
    \]

    Let $\eta\coloneqq2^{-(n+2)}\in\Q$. Using a loop over
    nonnegative integers, find $m$ such that
    \[
       C2^{-m}<\eta.
       \tag{3}
    \]
    Next find $d\ge0$ such that
    \[
       C\frac{B^{d+1}}{(d+1)!}<\eta.
       \tag{4}
    \]
    Both comparisons involve rationals. The first search terminates
    because $2^{-m}\to0$; the second terminates because for
    $k\ge2B$ the ratio of consecutive quantities
    $B^{k+1}/(k+1)!$ and $B^k/k!$ is $B/(k+1)<\tfrac12$.
    One can maintain the rational quantity $CB^k/k!$ inductively,
    multiplying it by $B/(k+1)$ on each iteration; there is no need to
    evaluate a real exponential during the search.

    Finally query $q\coloneqq\alpha(m)$ and return
    \[
         y\coloneqq P_d(q)=\sum_{j=0}^{d}\frac{q^j}{j!}\in\Q.
    \]
    This finite sum uses rational multiplication, division, and addition.
    By (1)--(4),
    \begin{align*}
       |y-\exp(x)|
       &\le |P_d(q)-\exp(q)|+|\exp(q)-\exp(x)|\\
       &<\eta+C2^{-m}<2\eta=2^{-(n+1)}<2^{-n}.
    \end{align*}
    All searches terminate, and every numerical operation in the
    algorithm is rational arithmetic or comparison. Because $x$ and its
    name $\alpha$ were arbitrary, this proves computability of
    $\exp\colon\R\to\R$.

    \item The accompanying \code{p3.ipynb} implements the preceding
    algorithm as \code{creal\_exp(x)}. It first obtains the integer
    bound $B$ from \code{x(0)}; it then chooses $m$ using (3),
    chooses $d$ using (4), and evaluates $P_d(\alpha(m))$ using an exact
    rational recurrence for its terms. The returned \code{CReal}
    computes this approximation afresh at each requested precision $n$.
    Thus its correctness and termination follow from part~(a),
    for every valid input approximation sequence.
\end{problemsubparts}
\end{solution}

\begin{problem}{4}
In this problem, you should prove claims using only the definition of
computable functions. Do not use any other results, including those
covered in class and in Problem~5.
\begin{problemsubparts}
    \item Prove that $f\colon\R\to\R$ defined by
    $f(x)\coloneqq(\text{if }x\ge0\text{ then }1\text{ else }0)$
    is not computable.
    \item Prove that $g\colon[0, \tfrac12]\to[0,1]$ defined by
    $g(x)\coloneqq x+\tfrac13$ is not computable under Turing's definition:
    A function $h\colon[0, \tfrac12]\to[0,1]$ is called
    \emph{computable in Turing's sense} if there exists an algorithm
    $\A$ that takes $(b_i)_{i\in\Nzero}\subseteq\{0,1\}$ and
    $j\in\Nzero$, and returns $0$ or $1$ such that
    \[
       h((0.b_0b_1\ldots)_2)=(0.c_0c_1\ldots)_2,
       \qquad c_j\coloneqq\A((b_i)_{i\in\Nzero}, j).
    \]
    % The source prints f instead of h on the left and has a misplaced
    % parenthesis. The equation above corrects these evident typos.
    This result shows that Turing's definition of computable functions
    is not equivalent to, and much more restrictive than, Borel's
    definition of computable functions.
\end{problemsubparts}
\end{problem}

\begin{solution}
\begin{problemsubparts}
    \item Suppose for contradiction that an oracle algorithm $\A$
    computes $f$ in the sense of the definition of a computable function.
    Give it precision $n=2$ and the name $\alpha(i)=0$ for every
    $i\in\Nzero$; this names $x=0$.
    The computation halts and returns some $q\in\Q$ with
    \[
       |q-f(0)|=|q-1|<\frac14,
       \qquad\text{so }q>\frac34.
    \]
    Only finitely many indices of $\alpha$ were queried. Let $N$ be
    at least all of them (take $N=0$ if none were queried), and put
    $t\coloneqq2^{-(N+2)}>0$.
    Define a second rational name by
    \[
       \beta(i)\coloneqq
       \begin{cases}
           0,&i\le N, \\
          -t,&i>N.
       \end{cases}
    \]
    This is a name of $-t$: for $i\le N$,
    $|\beta(i)-(-t)|=t<2^{-i}$, and for $i>N$ the error is zero.
    The algorithm receives exactly the same answers to all oracle
    queries made during its run on $\alpha$, so its deterministic
    execution on $\beta$ also halts with output $q$.
    Yet $f(-t)=0$, which requires $|q|<\tfrac14$, contradicting
    $q>\tfrac34$. Hence $f$ is not computable.

    \item Suppose that an algorithm $\A$ computes $g$ in the stated
    binary-digit sense. Consider the real number
    $x_0\coloneqq\tfrac16$. Its binary expansion is unique, since
    $\tfrac16$ is not dyadic, and
    \[
        g(x_0)=\frac12.
    \]
    Run $\A$ on the binary expansion of $x_0$ with output index
    $j=0$. It must halt after inspecting only finitely many input
    digits and output a bit $c_0\in\{0,1\}$.
    Let $N$ be at least all inspected indices (take $N=0$ if none
    were inspected). Because $x_0$ is not a dyadic rational, it lies
    strictly inside the dyadic interval specified by its first
    $N+1$ binary digits. Consequently there are numbers
    $x^-<x_0<x^+$ in $[0, \tfrac12]$ that share those digits with
    $x_0$. We may choose these inputs to be non-dyadic so that their
    binary expansions are unique.

    On each of these inputs, $\A$ must still return $c_0$, since
    the finite set of queried input digits has not changed.
    However,
    \[
        g(x^-)<\frac12<g(x^+).
    \]
    Every binary expansion of a number strictly below $\tfrac12$
    has first fractional digit $0$, whereas every binary expansion
    of a number strictly above $\tfrac12$ (and at most $1$) has
    first fractional digit $1$.
    Thus correctness on $x^-$ requires $c_0=0$, while correctness on
    $x^+$ requires $c_0=1$, a contradiction.
    This argument does not depend on which of the two binary expansions
    of $\tfrac12$ the algorithm chooses at $x_0$.
    Therefore $g$ is not Turing-computable as a function in the stated
    sense.
\end{problemsubparts}
\end{solution}

\begin{problem}{5}
Prove that if $f\colon\R\to\R$ is computable (under Borel's definition),
then $f$ is continuous.

\begin{hint}
To prove that $f$ is continuous, it suffices to prove the
following: for all $a\in\R$ and $n\in\N$, there exists $m\in\N$ such that
$|x-a|<2^{-m}$ implies $|f(x)-f(a)|<2^{-n}$ for all $x\in\R$.
\end{hint}
\end{problem}

\begin{solution}
Let $\A$ be an oracle algorithm computing $f$. Fix an arbitrary
$a\in\R$ and $n\in\N$. For every $i\in\Nzero$, choose a rational
number $\alpha(i)$ with
\[
       |\alpha(i)-a|<2^{-(i+2)}.
\]
Such rationals exist by density of $\Q$ in $\R$, even if $a$ is not
computable. In particular $\alpha$ is a valid rational name of $a$.
Run $\A$ on this name at precision $n+2$. By computability it halts
with a rational result $q$ satisfying
\[
      |q-f(a)|<2^{-(n+2)}.
\]
Let $I\subseteq\Nzero$ be the finite set of oracle indices queried
in this run. For each $i\in I$, the quantity
$2^{-i}-|\alpha(i)-a|$ is strictly positive. Consequently, there
exists $\delta>0$ such that
\[
     |x-a|<\delta
     \quad\Longrightarrow\quad
     |\alpha(i)-x|<2^{-i}\quad(i\in I).
\]
Explicitly, take
$\delta\coloneqq\min_{i\in I}(2^{-i}-|\alpha(i)-a|)$
when $I\ne\varnothing$, and $\delta\coloneqq1$ otherwise.

For any $x\in\R$ with $|x-a|<\delta$, construct a rational name
$\beta$ of $x$ by setting $\beta(i)=\alpha(i)$ for $i\in I$
and choosing, for each $i\notin I$, any rational $\beta(i)$
within $2^{-i}$ of $x$.
This construction need not be uniformly computable: the definition
requires $\A$ to work for every valid name.
The names $\alpha$ and $\beta$ give identical answers to every
query made in the first computation, so $\A(\beta, n+2)=q$.
It follows that
\begin{align*}
 |f(x)-f(a)|
  &\le |f(x)-q|+|q-f(a)|\\
  &<2^{-(n+2)}+2^{-(n+2)}=2^{-(n+1)}<2^{-n}.
\end{align*}
Choose $m\in\N$ with $2^{-m}<\delta$. Then
$|x-a|<2^{-m}$ implies the required estimate for every
$x\in\R$. Since $a$ and $n$ were arbitrary, $f$ is continuous
on $\R$.
\end{solution}

\end{document}
